\section*{Sets and Indices}

\begin{itemize}
    \item $P$: set of products, indexed by $i$.
    \item In this instance,
    \[
    P=\{\text{Shirt},\text{Short-sleeve},\text{Casual Cloth}\}.
    \]
\end{itemize}

\section*{Parameters}

\begin{itemize}
    \item $a_i^{\text{lab}}$: labor required per unit of product $i$.
    \item $a_i^{\text{mat}}$: material required per unit of product $i$.
    \item $p_i$: selling price per unit of product $i$.
    \item $c_i^{\text{var}}$: variable cost per unit of product $i$.
    \item $u_i = p_i - c_i^{\text{var}}$: unit contribution profit of product $i$ (code: \texttt{unit\_profit[i]}).
    \item $L$: available labor per week (code: \texttt{available\_labor}).
    \item $M$: available material per week (code: \texttt{available\_material}).
    \item $F_i$: weekly fixed cost associated with the equipment for product $i$.
    \item $F^{\text{tot}}=\sum_{i\in P} F_i$ (code: \texttt{total\_fixed}).
\end{itemize}

For this instance, the data are:
\[
L=1500,\qquad M=1600,
\]
and
\[
\begin{array}{c|cccccc}
i & a_i^{\text{lab}} & a_i^{\text{mat}} & p_i & c_i^{\text{var}} & u_i & F_i\\
\hline
\text{Shirt} & 3 & 4 & 120 & 60 & 60 & 2000\\
\text{Short-sleeve} & 2 & 3 & 80 & 40 & 40 & 1500\\
\text{Casual Cloth} & 6 & 6 & 180 & 80 & 100 & 1000
\end{array}
\]
so that
\[
F^{\text{tot}}=2000+1500+1000=4500.
\]

\section*{Decision Variables}

\begin{itemize}
    \item $x_i$ (code: \texttt{x[i]}): integer number of units of product $i$ to produce during the week.
\end{itemize}

Variable domain:
\[
x_i \in \mathbb{Z}_{\ge 0}\qquad \forall i\in P.
\]

\section*{Objective}

Maximize weekly profit:
\[
\max \sum_{i\in P} u_i x_i - F^{\text{tot}}.
\]

For this instance, equivalently,
\[
\max\; 60x_{\text{Shirt}} + 40x_{\text{Short-sleeve}} + 100x_{\text{Casual Cloth}} - 4500.
\]

\section*{Constraints}

Labor availability:
\[
\sum_{i\in P} a_i^{\text{lab}} x_i \le L.
\]

Material availability:
\[
\sum_{i\in P} a_i^{\text{mat}} x_i \le M.
\]

For this instance:
\[
3x_{\text{Shirt}} + 2x_{\text{Short-sleeve}} + 6x_{\text{Casual Cloth}} \le 1500,
\]
\[
4x_{\text{Shirt}} + 3x_{\text{Short-sleeve}} + 6x_{\text{Casual Cloth}} \le 1600.
\]

\section*{Notes on Implementation}

\begin{itemize}
    \item The implemented model in \texttt{current\_heuristic.py} is a linear integer program solved with Gurobi through the PuLP-compatible interface.
    \item Production quantities are required to be nonnegative integers; there are no binary setup variables.
    \item All fixed costs are subtracted unconditionally in the objective, regardless of whether any units of the corresponding product are produced. This matches the code comment: the equipment already exists, so the model always deducts
    \[
    F^{\text{tot}}=\sum_{i\in P} F_i.
    \]
    \item Therefore, the formulation does \emph{not} model product-specific activation decisions or avoid fixed costs when a product is not produced.
\end{itemize}
